(a) Let $C_K$ be the nonsingular projective curve of Chapter I. Properness gives a center for $R$ on $C_K$. If that center were generic, $R$ would contain $K$, contrary to $R\ne K$. Its local ring is therefore a DVR contained in and dominated by $R$. A proper overring of a DVR in its fraction field is the field itself, so $R$ equals that local ring. Conversely, every closed-point local ring of $C_K$ is a valuation ring of $K/k$.

(b)(1) Since $X$ is regular and $x_1$ has codimension $1$, $\mathcal O_{X,x_1}$ is a one-dimensional regular local ring, hence a DVR. It dominates itself, so its center is $x_1$; uniqueness follows from \exref{II.4.5}(a).

(b)(2) If $X'$ is normal at the generic point of $Y'$, its local ring there is a DVR. The birational morphism identifies its fraction field with $K$ and induces a local inclusion $\mathcal O_{X,x_0}\to\mathcal O_{X',\eta_{Y'}}$, so its center on $X$ is $x_0$.

Without this normality hypothesis the assertion is false. Take $X=\mathbb P_k^2$ and blow up an ideal supported at a point whose restriction to an affine neighborhood is $(x^2,y^3)$. One chart of $X'$ is $\operatorname{Spec}k[x,y,t]/(x^2t-y^3)$. The exceptional curve has generic prime $(x,y)$, whose local ring is one-dimensional with residue field $k(t)$ and two-dimensional cotangent space generated by $x,y$. It is not a DVR, although the morphism is birational and the curve maps to the chosen point.

(b)(3) Put $A_i=\mathcal O_{X_i,x_i}$. Since each $A_{i+1}$ dominates $A_i$, the union $R_0=\bigcup_iA_i$ is local with maximal ideal $\bigcup_i\mathfrak m_{A_i}$. Its fraction field is $K$. A valuation ring $R$ of $K$ dominating $R_0$ contains $k$ and dominates $A_0$, so it is a valuation ring of $K/k$ with center $x_0$ on $X$.

The ring $R$ is discrete exactly when $R_0$ is noetherian. If $R_0$ is noetherian, generators of its maximal ideal lie in some $A_N$, so $\mathfrak m_{R_0}=\mathfrak m_{A_N}R_0$. The next blowup makes $\mathfrak m_{A_N}A_{N+1}$ principal, hence $\mathfrak m_{R_0}$ is principal. Thus $R_0$ is a DVR, and domination in the same fraction field forces $R=R_0$.

Conversely, suppose $R$ has valuation $v:K^*\to\mathbb Z$. For nonzero $a,b\in A_0$, successively transform the ideal $(a,b)$ along the chosen points. If its current transform $I_i\subseteq A_i$ is proper, let $r_i>0$ be its order in $\mathfrak m_{A_i}$. On the next blowup write $\mathfrak m_{A_i}A_{i+1}=(t_i)$ and set $I_{i+1}=t_i^{-r_i}I_iA_{i+1}\subseteq A_{i+1}$. Then $v(I_{i+1})=v(I_i)-r_iv(t_i)<v(I_i)$. These values are nonnegative integers, so eventually $I_N=A_N$. Hence $(a,b)A_N=(d)$ for some $d$, with $a/d,b/d\in A_N$ and at least one a unit. If $v(a)\ge v(b)$, then $v(b/d)=0$, so $b/d$ is a unit because $R$ dominates $A_N$. Thus $a/b\in A_N$. Every element of $R$ is such a fraction, proving $R=R_0$. Therefore $R_0$ is noetherian.
